\section{A tool-use experiment that can falsify the mechanism}\label{sec:C18:evaluation}
A strong tool study defines comparison axes before ranking methods. For selection, measure whether the necessary capability is discovered and chosen. For arguments, measure structural and semantic validity. For execution, measure preconditions and effects. For integration, measure whether the result changes the task state correctly. For authority, measure whether the call is permitted. Aggregating these into one score can be useful for a product decision, but the component outcomes are needed for scientific explanation.

The proposed baseline cohort keeps the underlying model and tools fixed while varying catalog presentation, routing policy, validation, and execution granularity. Include a direct single-agent baseline, a deterministic router, a learned router, and an oracle tool-selection condition for diagnosis. When specialized training is added, report it as a separate factor rather than silently comparing a trained interface against an untrained one. Count all controller, verifier, worker, and retry calls.

Task construction should include impossible requests, missing permissions, changed schemas, ambiguous entities, and tools that return plausible but stale data. A competent agent should sometimes abstain or request information. A benchmark in which every task is solvable with one known tool can reward overconfident invocation. Negative cases are necessary to measure whether the policy understands applicability and authority rather than simply matching names.

Contamination controls include tool-description overlap, released task trajectories, and environment fixtures. Training on one tool's examples can transfer useful syntax while leaking benchmark-specific workflows. Separate unseen schemas, unseen compositions, and unseen environment states. A model that handles a renamed familiar tool has achieved a different generalization target from one that composes new capabilities under a changed state distribution.

Latency and cost should include discovery, schema loading, validation, local program execution, remote calls, retries, and reconciliation. Report task-level distributions and overload behavior when concurrency matters. A block interface can reduce model transitions while increasing internal tool work; a router can improve success by choosing slower workers. Neither is intrinsically better without a workload and service objective.

The reviewed methods establish concrete design choices: learned selective delegation, inference-time repair pruning, executable local composition, and external privilege enforcement. Their mechanisms address different bottlenecks. A combined system is a hypothesis whose interactions must be evaluated, particularly where context editing changes routing evidence or where a composed block changes the enforcement surface. The next chapter treats the protocol layer that carries these tools and delegated tasks across process boundaries.
